% CVPR 2026 Paper Template; see https://github.com/cvpr-org/author-kit

\documentclass[10pt,twocolumn,letterpaper]{article}

%%%%%%%%% PAPER TYPE
\usepackage{cvpr}              % CAMERA-READY
% \usepackage[review]{cvpr}      % REVIEW
% \usepackage[pagenumbers]{cvpr} % arXiv version

% -------------------- Basic packages --------------------
\usepackage[utf8]{inputenc}
\usepackage[T1]{fontenc}
\usepackage{times}
\usepackage{epsfig}
\usepackage{graphicx}
\usepackage{amsmath,amssymb}
\usepackage{booktabs}
\usepackage{multirow}
\usepackage{array}
\usepackage{enumitem}
\usepackage{float}
\usepackage{url}
\usepackage{xcolor}
\usepackage{caption}


% -------------------- Hyperref --------------------
\definecolor{cvprblue}{rgb}{0.21,0.49,0.74}
\usepackage[pagebackref,breaklinks,colorlinks,citecolor=cvprblue,linkcolor=cvprblue,urlcolor=cvprblue]{hyperref}

%%%%%%%%% PAPER ID
\def\paperID{hit\_zzy}
\def\confName{CVPR}
\def\confYear{2026}

%%%%%%%%% TITLE
\title{NTIRE 2026 Remote Sensing Infrared Image Super-Resolution Challenge Factsheet\\
Team hit\_zzy}

%%%%%%%%% AUTHORS
\author{
Zengyuan Zuo\\
Harbin Institute of Technology\\
{\tt\small 3565741165@qq.com} \and Junjun Jiang\\
Harbin Institute of Technology\\
{\tt\small jiangjunjun@hit.edu.cn}
}

\begin{document}

\maketitle

\section{Factsheet Information}

\subsection{Team details}

\begin{itemize}
    \item \textbf{Team name:} hit\_zzy
    \item \textbf{Team leader name:} Zengyuan Zuo
    \item \textbf{Team leader email:} 3565741165@qq.com
    \item \textbf{Rest of the team members:} Junjun Jiang
    \item \textbf{Team website URL (if any):} N/A
    \item \textbf{Affiliation:} Harbin Institute of Technology
    \item \textbf{Affiliation of the team and/or team members with NTIRE 2026 sponsors:} None
    \item \textbf{User names and entries on the NTIRE 2026 CodaLab competitions:} hit\_zzy
    \item \textbf{Best scoring entries of the team during the testing phase:} 54.02
    \item \textbf{Link to the codes/executables of the solution(s):} \url{https://github.com/Harbinzzy/NTIRE2026_infraredSR}
\end{itemize}

\subsection{Method details}

\subsubsection{General method description}

Our solution IRSRMambaPlus is built upon the IRSRMamba~\cite{huang2025irsrmamba} framework, a Mamba-based infrared image super-resolution architecture that combines long-range dependency modeling and multi-scale feature extraction. Based on this baseline, we adapt the method to the remote sensing infrared super-resolution challenge by introducing two additional guidance priors into the input space and by redesigning the training pipeline.

Specifically, instead of only using the original three-channel low-resolution image as input, we concatenate two additional guidance maps~\cite{zou2026toward}, namely an edge map and a heat map, resulting in a five-channel input representation:
\[
\mathbf{x}_{in} = [\mathbf{x}_{RGB}, \mathbf{x}_{edge}, \mathbf{x}_{heat}] \in \mathbb{R}^{H \times W \times 5}.
\]
The edge map is used to provide explicit structural and boundary information, while the heat map is used to encode smoothed intensity distribution cues. This modification is lightweight, but effective for remote sensing infrared images, which often exhibit weak textures, low contrast, and blurred boundaries.

During inference, we adopt self-ensemble to further improve reconstruction quality. 
Specifically, eight geometric transformations (four rotations and their corresponding horizontal flips) are applied to the input image. 
Each transformed input is processed independently by the network, and the inverse transformations are applied to the outputs before averaging them to obtain the final result.

\subsubsection{Training strategy}

In addition to the input redesign, we employ a three-stage curriculum training strategy. 
In the first stage, the model is trained on the training set with $\times2$ supervision to learn a stable coarse reconstruction mapping. 
In the second stage, the model is further trained on the training set with $\times4$ supervision to adapt to the target upscaling factor. 
In the third stage, the model is refined using alternating training between the training set $\times4$ data and the validation set $\times2$ data, where the auxiliary branch is assigned to improve robustness.

The model is optimized using Adam with an initial learning rate of $2\times10^{-4}$ and betas $(0.9, 0.99)$ without weight decay. 
A MultiStepRestartLR scheduler is used with milestones at 70k, 110k, and 170k iterations and a decay factor of 0.5, together with restart points at 0, 30k, and 140k iterations. 
The total training length is 200k iterations with a warm-up of 500 iterations. 
The default batch size is 12 per GPU and is increased to 48 in the second and third stages. 
Training patches of size $64\times64$ are randomly cropped, and random horizontal flipping and rotation are applied for data augmentation.

\subsubsection{Loss function}

The training objective combines a pixel-domain reconstruction loss with an additional SSM-related loss. 
The SSM loss weight is set to 1.0 and the pixel loss weight to 0.2. 
Mean reduction and mean aggregation are used during optimization.

\section{Other details}

\begin{itemize}
    \item \textbf{Planned submission of a solution description paper at the NTIRE 2026 workshop:} No
    \item \textbf{General comments and impressions of the NTIRE 2026 challenge:} The challenge is well designed and practically meaningful. Remote sensing infrared super-resolution is a difficult but important problem, and the benchmark encourages methods that improve both structural fidelity and restoration quality.
    \item \textbf{What do you expect from a new challenge in remote sensing image restoration, enhancement and manipulation?} We hope future challenges can include more realistic degradation settings, stronger cross-domain evaluation, and more task-oriented assessment protocols.
    \item \textbf{Other comments:} No
\end{itemize}

\bibliographystyle{ieeenat_fullname}
\bibliography{main}

\end{document}